\rSec3[demo.holder]{Class template holder}

\pnum
Subclause \iref{demo.holder} describes the class template \tcode{holder}.

\begin{codeblock}
namespace std {
  template<class E>
  class @\libglobal{holder}@ {};
}
\end{codeblock}

\pnum
A program that instantiates the definition of \tcode{holder} for a non-object type is ill-formed.


\rSec3[demo.plain]{Class template plain}

\pnum
Subclause \iref{demo.plain} describes the class template \tcode{plain}.

\pnum
It holds nothing.

\begin{codeblock}
namespace std {
  template<class E>
  class @\libglobal{plain}@ {};
}
\end{codeblock}

\pnum
A program that instantiates \tcode{plain<E>} is ill-formed unless \tcode{is_object_like<E>} is \tcode{true}.

\pnum
\remarks
A \tcode{plain} is trivially copyable.

\rSec3[demo.late]{Class template late}

\begin{codeblock}
namespace std {
  template<class E>
  class @\libglobal{late}@ {};
}
\end{codeblock}

\pnum
Subclause \iref{demo.late} describes the class template \tcode{late}.

\rSec3[demo.routed]{Class template routed}

\begin{codeblock}
namespace std {
  template<class E>
  class @\libglobal{routed}@ {};
}
\end{codeblock}

\rSec4[demo.routed.general]{General}

\pnum
Subclause \iref{demo.routed} describes the class template \tcode{routed}.

\pnum
\remarks
A \tcode{routed} is described in its general subclause.
